\section{Simulation}\label{sec:simulation}

\subsection{Design}

Data were simulated from a joint model with a random intercept and slope,
longitudinal mean $\beta_0 + \beta_{\mathrm{time}} t +
\beta_{\mathrm{age}}\mathrm{age}_i + \beta_{\mathrm{sex}}\mathrm{sex}_i$
with $\beta_0 = 2$ and $\beta_{\mathrm{time}} = 0.5$, $\mathrm{age}_i \sim \N(0, 10^2)$ and $\mathrm{sex}_i \sim
\mathrm{Bernoulli}(0.5)$, $(\sigma_0, \sigma_1) = (0.8, 0.3)$, $\sigma_e = 0.3$, and a
Weibull hazard (shape 1.3) with $\alpha = 0.5$ and a treatment effect
$\gamma = -0.3$, calibrated to an event rate of about one half. Visits
were scheduled at regular intervals up to $t = 10$ and stopped at the event
or censoring time. All fits used the spline baseline hazard of
Section~\ref{sec:setup}, so the fitted model does not contain the true
hazard exactly. Each design was simulated with three seeds, each
fitted unrotated and with the rotation and its dense block. ESS is NumPyro's (version 0.15) multi-chain
autocorrelation estimator, truncated by Geyer's initial monotone sequence
and without rank normalization \citep{geyer1992,gelman2013}, and $\hat R$
is its split-chain version. ESS per second divides by the wall time of the
whole run, warm-up and compilation included; ESS per gradient divides by
the number of leapfrog steps in the sampling phase. Ratios are rotated over
unrotated, paired by seed, with geometric means over seeds.

The core grid crossed about 9 against about 17 visits per subject with
$\rho \in \{0, 0.3, 0.6, 0.9\}$ at $N = 300$, and added $N \in \{150,
600\}$ at 9 visits and $\rho = 0.3$: ten designs, 2 chains of 500 warm-up
and 1000 sampling iterations.

\subsection{Efficiency}

\begin{table}[t]
\centering
\caption{Core grid: ESS per second for the intercept and the time slope
(means over 3 seeds), and the rotated/unrotated ratio (geometric mean of
the paired per-seed ratios).}
\label{tab:core}
\small
\begin{tabular}{rrr rrr rrr r}
\toprule
& & & \multicolumn{3}{c}{intercept} & \multicolumn{3}{c}{time slope} & $\alpha$ \\
\cmidrule(lr){4-6}\cmidrule(lr){7-9}\cmidrule(lr){10-10}
visits & $\rho$ & $N$ & unrot. & rotated & ratio & unrot. & rotated & ratio & ratio \\
\midrule
9  & 0.0 & 300 & 25.9 & 354.4 & 14.2 &  9.1 & 304.9 & 36.4 & 0.85 \\
17 & 0.0 & 300 & 12.1 & 229.4 & 19.3 &  4.0 & 234.9 & 59.2 & 1.32 \\
9  & 0.3 & 300 & 30.0 & 253.3 &  8.6 & 19.0 & 251.1 & 13.4 & 0.72 \\
17 & 0.3 & 300 & 11.3 & 248.5 & 22.2 &  6.4 & 262.5 & 41.9 & 1.33 \\
9  & 0.6 & 300 & 26.1 & 245.7 &  9.4 & 29.9 & 229.3 &  7.7 & 1.00 \\
17 & 0.6 & 300 &  9.6 & 231.6 & 24.0 & 10.0 & 251.2 & 24.9 & 1.59 \\
9  & 0.9 & 300 & 20.4 & 220.7 & 11.1 & 20.3 & 250.3 & 12.8 & 0.90 \\
17 & 0.9 & 300 &  8.9 & 218.9 & 24.9 &  8.6 & 217.7 & 25.6 & 1.59 \\
9  & 0.3 & 150 & 41.8 & 323.7 &  7.6 & 32.6 & 420.4 & 13.1 & 0.76 \\
9  & 0.3 & 600 & 15.9 & 132.0 &  8.2 &  8.8 & 151.7 & 17.7 & 0.77 \\
\bottomrule
\end{tabular}
\end{table}

Over all 30 paired fits of Table~\ref{tab:core}, the geometric-mean gain
in ESS per second was 13.5 for the intercept (per seed 4.5--31.6) and 21.1
for the time slope (5.1--70.8). Wall time was essentially unchanged (ratio
1.03, range 0.83--1.10) and the rotated fit took 0.89 times as many
leapfrog steps per iteration, so per gradient the gains were 15.6 and
24.3. Almost all of the gain is in ESS per draw: a median of about 1.7 rotated
against 0.1 unrotated. There were no divergences, and the worst $\hat R$ on a
reported coefficient was 1.002 rotated and 1.044 unrotated. Posterior
means agreed to a median of 0.027 posterior standard deviations (largest
paired difference 2.7 Monte Carlo standard errors over 90 pairs), and
posterior standard deviations to within Monte Carlo error (ratios 1.01 for
the intercept and 0.98 for the time slope).

The pattern follows Proposition~\ref{prop:ridge}. The gain is largest with
dense visits, where each subject carries the most information about its
own random effects (19--25 times for the intercept, 25--59 for the slope),
and does not grow with $N$ (7.6 at $N = 150$, 8.2 at $N = 600$).
Figure~\ref{fig:gain} shows all three grids against the per-subject
information $\bar n\,\sigma_0^2/\sigma_e^2$. The association parameter,
which the degeneracy does not involve, split by design: 0.72--1.00 in the
six designs with about 9 visits and 1.32--1.59 in the four with about 17
(geometric mean 1.04). We have no established explanation; on real data
(Section~\ref{sec:realdata}) the effect on $\alpha$ tracked whether the
degeneracy had been setting NUTS's trajectory length, a plausible reading
here too that was not measured.

\begin{figure}[t]
\centering
\includegraphics[width=\textwidth]{fig2_gain.pdf}
\caption{ESS per second of the reported coefficient, rotated over unrotated
fit on the same simulated data, against the per-subject information
$\bar n\,\sigma_0^2/\sigma_e^2$ ($\sigma_0 = 0.8$). Large markers:
geometric mean over three seeds in one design; small: single seeds. Core
and stress grids have a random intercept and slope ($q = 2$); the $q = 1$
grid has a random intercept only. In (b), open triangles are the $q = 1$
time slope, which has no random effect and so nothing to rotate. The $q =
1$ intercept ratios are at the standard run length, where most unrotated
fits had not converged; Section~\ref{sec:q1} compares converged fits.}
\label{fig:gain}
\end{figure}

\subsection{Random intercept only}\label{sec:q1}

The same simulator with $\sigma_1 = 0$, fitted with a random intercept
only, gave seven designs (about 9, 17 and 35 visits, $N \in \{100, 300\}$,
and one with $\sigma_e = 0.8$). The intercept gained 6.5--282 times and the
coefficient of age 6.6--207 times, but at this run length 17 of the 21
unrotated fits had some $\hat R$ above 1.05, while every rotated fit had
$\hat R \leq 1.004$. For a comparison between converged fits we reran three
designs ($N = 300$; about 9, 17 and 35 visits), extending the unrotated
fit to up to 32 times the standard length, with a criterion fixed in
advance: $\hat R \leq 1.01$ on every monitored parameter and ESS $\geq 400$
for the intercept, age, the time slope and $\alpha$. Every rotated fit met
it at the standard length. The unrotated fit met it first at 16 times the
standard length with 9 visits, at 16 and 32 times with 17 visits, and not
by 32 times with 35 visits. Its ESS per draw for the intercept did not
change with run length, and the rotated fit's was of order 100 times
larger, so the advantage is a property of the stationary chains, not of
short runs. In wall time to the criterion the unrotated fit needed 3.5 and
6.5--6.6 times as long in the first two designs, and more than 10--14 times
in the third (\suppref{sec:longrun}).

\subsection{Stress grid}\label{sec:stress}

Section~\ref{sec:residual} names two regimes where the rotation should be
weakest: few subjects each carrying a lot of information, and strongly
correlated random effects. A stress grid of six designs targeted both
(about 35 visits with $N \in \{50, 100\}$ and $\rho \in \{0.3, 0.95\}$,
plus two designs with $\rho = 0.95$), with the criterion, fixed in
advance, that the rotation must not fall below the unrotated fit on the
intercept or time slope in any design. It did not: the lowest design gained
5.5 times and the lowest single seed 3.2 times, and per gradient the
geometric-mean gains were 16.7 for the intercept and 22.0 for the time
slope (\suppref{sec:stressgrid}). The design with the fewest subjects and
the most information per subject gained least, as the fixed-metric
residual predicts, and strong correlation did not reduce the gains. The
predicted weakness is visible only against a variant that removes the
absorbable directions and restores the reported coefficients with an
exact correction (Section~\ref{sec:related}), which has no fixed-metric
residual: in the 35-visit designs it was 1.4--2 times more efficient than
the rotation, and about equal with 9 visits.

\subsection{Calibration}\label{sec:calibration}

Proposition~\ref{prop:exact} says the rotation cannot change the
posterior, so a calibration study checks the implementation. With $N =
300$ and 100 replicates, none excluded, coverage of 95\% intervals agreed
between the two fits within 0.01 and the ratio of mean posterior standard
deviation to the standard deviation of posterior means within 0.015 for
$\beta_0$, $\beta_{\mathrm{time}}$, $\alpha$ and $\sigma_0$
(\suppref{sec:calib}). Both fits showed an upward bias of about 3 Monte
Carlo standard errors (2.8\%) in $\alpha$, which fell to 1.1\% at $N =
1200$, consistent with ordinary finite-sample bias.
